% Einheit 1 von 4 – Einheitskreis und Bogenmaß (bank/trigonometrische-funktionen/e1.jsonl, zusammenbau v0.1)
%% TODO Zweigzeile Teil 1: was hier gelernt wird (Satz in Schülersprache aus dem Einheitstitel) – Einheit 1
\einheitenkopf[e1]{Einheit 1 von 4 · Einheitskreis und Bogenmaß}
\zweigzeile{ab Kl. 10 · keine P10-Aufgabe}
\setcounter{aufgabe}{13}

%% TODO Titel: Ich-kann-Satz fehlt in der Bank (Platzhalter: Kettenname) – Nr. 14
%% TODO Anweisung: ein Satz über der ersten Teilaufgabe fehlt in der Bank – Nr. 14
\begin{aufgabe}{Einheitskreis und Bogenmaß}
\begin{teile}
\teil Grad oder Bogenmaß? Die Winkelangabe lautet 2,4. Kreuze das Maß und den passenden Modus an.\\ \kreuz{Gradmaß}\kreuz{Bogenmaß}\\ \kreuz{DEG}\kreuz{RAD}
\end{teile}
\swa{Zeichne den Punkt P zum Winkel 40° auf dem Einheitskreis ein.}{\begin{ksys}[xmin=-1.2,xmax=1.2,ymin=-1.2,ymax=1.2,xstep=0.2,ystep=0.2]\funktionab{sqrt(abs(1-\x^2))}{}{-1}{1}\funktionab{-sqrt(abs(1-\x^2))}{}{-1}{1}\end{ksys}}
\swa{Zeichne den Punkt P zum Winkel 130° auf dem Einheitskreis ein.}{\begin{ksys}[xmin=-1.2,xmax=1.2,ymin=-1.2,ymax=1.2,xstep=0.2,ystep=0.2]\funktionab{sqrt(abs(1-\x^2))}{}{-1}{1}\funktionab{-sqrt(abs(1-\x^2))}{}{-1}{1}\end{ksys}}
\swa{Zeichne den Punkt P zum Winkel 215° auf dem Einheitskreis ein.}{\begin{ksys}[xmin=-1.2,xmax=1.2,ymin=-1.2,ymax=1.2,xstep=0.2,ystep=0.2]\funktionab{sqrt(abs(1-\x^2))}{}{-1}{1}\funktionab{-sqrt(abs(1-\x^2))}{}{-1}{1}\end{ksys}}
\swa{Zeichne den Punkt P zum Winkel 320° auf dem Einheitskreis ein.}{\begin{ksys}[xmin=-1.2,xmax=1.2,ymin=-1.2,ymax=1.2,xstep=0.2,ystep=0.2]\funktionab{sqrt(abs(1-\x^2))}{}{-1}{1}\funktionab{-sqrt(abs(1-\x^2))}{}{-1}{1}\end{ksys}}
\swa{Zeichne den Punkt P zum Winkel 70° auf dem Einheitskreis ein.}{\begin{ksys}[xmin=-1.2,xmax=1.2,ymin=-1.2,ymax=1.2,xstep=0.2,ystep=0.2]\funktionab{sqrt(abs(1-\x^2))}{}{-1}{1}\funktionab{-sqrt(abs(1-\x^2))}{}{-1}{1}\end{ksys}}
\swfrage{Was bleibt gleich, was ändert sich?}
\swz{Miss den Winkel $\alpha$ zum Punkt P.}{\einheitskreis[ohne]{100}}
\swz{Lies sin $\alpha$ und cos $\alpha$ für den Punkt P ab (eine Stelle).}{\begin{ksys}[xmin=-1.2,xmax=1.2,ymin=-1.2,ymax=1.2,xstep=0.2,ystep=0.2,ablesen]\funktionab{sqrt(abs(1-\x^2))}{}{-1}{1}\funktionab{-sqrt(abs(1-\x^2))}{}{-1}{1}\punkt{0.6}{0.8}{P}\end{ksys}}
%% TODO Darstellung für schwach fehlt in der Bank (trigonometrische-funktionen-e1-k1-s4-v1; Abschnitt „Für schwache Schüler“)
\swz{P liegt ganz oben auf dem Einheitskreis – Winkel, Sinuswert und Kosinuswert?}{}
\swz{Am Einheitskreis wurde sin 50° ≈ 0,8 abgelesen. Prüfe mit dem Taschenrechner im Grad-Modus (zwei Stellen).}{\einheitskreis{50}}
%% TODO Darstellung für schwach fehlt in der Bank (trigonometrische-funktionen-e1-k1-s6-v1; Abschnitt „Für schwache Schüler“)
\swz{160° – welches Vorzeichen haben Sinuswert und Kosinuswert?}{}
%% TODO Darstellung für schwach fehlt in der Bank (trigonometrische-funktionen-e1-k1-s7-v1; Abschnitt „Für schwache Schüler“)
\swz{sin $\alpha = 0{,}6$ – welche beiden Winkel im ersten Vollkreis ab 0° passen? Runde auf eine Stelle.}{}
%% TODO Darstellung für schwach fehlt in der Bank (trigonometrische-funktionen-e1-k1-s8-v1; Abschnitt „Für schwache Schüler“)
\swz[3]{60° ist ein Sechstel des Vollwinkels. Begründe mit dem Umfang des Einheitskreises, dass das Bogenmaß $\frac{\pi}{3}$ ist.}{}
%% TODO Darstellung für schwach fehlt in der Bank (trigonometrische-funktionen-e1-k1-s9-v1; Abschnitt „Für schwache Schüler“)
\swz{40° ins Bogenmaß – als Vielfaches von $\pi$ und auf zwei Stellen?}{}
%% TODO Darstellung für schwach fehlt in der Bank (trigonometrische-funktionen-e1-k1-s10-v1; Abschnitt „Für schwache Schüler“)
\swz{$\frac{3}{4}\pi$ ins Gradmaß – wie viel Grad?}{}
\begin{teile}
\teil Welches Bogenmaß gehört zu 60°?\\ \kreuz{$\frac{\pi}{6}$}\kreuz{$\frac{\pi}{3}$}\kreuz{$\frac{\pi}{2}$}\kreuz{$\frac{2\pi}{3}$}
\end{teile}
%% TODO Darstellung für schwach fehlt in der Bank (trigonometrische-funktionen-e1-k1-s12-v1; Abschnitt „Für schwache Schüler“)
\swz{Der Rechner zeigt für sin 1 den Wert 0,017. Welcher Modus ist eingestellt, und was zeigt er im anderen Modus (zwei Stellen)?}{}
%% TODO Darstellung für schwach fehlt in der Bank (trigonometrische-funktionen-e1-k1-s13-v1; Abschnitt „Für schwache Schüler“)
\swz[3]{Eine Diagonale e wird mit dem Sinussatz berechnet: e = 5,3 · sin 128° : sin 31° (in cm). Berechne e auf zwei Stellen und zeige am Einheitskreis, warum es sin 128° gibt, obwohl kein rechtwinkliges Dreieck einen Winkel von 128° hat. \hfill (P10 2015 OS)}{}
\swz[3]{a) 105° – wie groß im Bogenmaß (zwei Stellen)? b) $\frac{4}{3}\pi$ – wie groß im Gradmaß? c) Erkläre am Einheitskreis, woher sin 105° kommt, und gib den Wert auf zwei Stellen an.}{}
\end{aufgabe}

%% TODO Titel: Ich-kann-Satz fehlt in der Bank (Platzhalter: Kettenname) – Nr. 15
%% TODO Anweisung: ein Satz über der ersten Teilaufgabe fehlt in der Bank – Nr. 15
\begin{aufgabe}{Einheitskreis und Bogenmaß – Fehler finden, Begründen}
\swz[3]{Lena berechnet sin 25° und erhält −0,13. Finde den Fehler und gib den richtigen Wert an.}{}
\swfrage{Begründe am Einheitskreis, warum sin $\alpha$ nie größer als 1 sein kann.}
\end{aufgabe}

%% TODO Merkkasten Einheit 1: 7 Zeilen, Vorgabe höchstens fünf (3.1)
\uebersichtskasten{Einheitskreis: Ein Kreis um den Ursprung mit dem Radius 1. Zu jedem Winkel gehört ein Punkt auf dem Kreis – seine Hochkoordinate ist der Sinuswert, seine Rechtskoordinate der Kosinuswert des Winkels. \\ 30°: P(0,87 | 0,5) \quad sin 30° = 0,5 \quad cos 30° ≈ 0,87 \\ Über den rechten Winkel hinaus: Der Punkt läuft weiter um den Kreis, deshalb gibt es Sinus- und Kosinuswerte für jeden Winkel. Rechts und oben sind die Koordinaten positiv, links und unten negativ. \\ sin 150° = 0,5 \quad cos 150° ≈ −0,87 \\ Bogenmaß: Statt in Grad misst man den Winkel durch die Länge seines Bogens im Einheitskreis. Der ganze Kreis hat den Umfang 2π, also gehört zum Vollwinkel 360° das Bogenmaß 2π. \\ 180° ≙ π \quad 90° ≙ π/2 \quad Bogenmaß = Gradzahl · π : 180 \quad Gradzahl = Bogenmaß · 180 : π \\ Am Taschenrechner: DEG rechnet in Grad, RAD im Bogenmaß – vor jeder Rechnung den Modus prüfen.}
